\section{Cryptographic Commitment Protocol}
\label{sec:commitment}

% Core formal contribution #3: the commitment-verification protocol.
% This section is ALIGNMENT-GENERAL. All verification is algorithmic
% against accepted benchmarks/protocols. Human judgment enters only
% at the governance layer (§4.2), not here.

% §5.1 Protocol overview
%   Three-phase protocol for verifying a claim against an accepted
%   benchmark:
%   Phase 1 (Commit): Agent a commits to claim C by publishing
%     Com(C, r) to the ledger. The commitment is binding and hiding.
%   Phase 2 (Verify): An algorithmic witness W⊥_alg executes the
%     accepted exercise protocol E for claim type C. The exercise
%     protocol is fully specified by the governance-accepted benchmark;
%     the witness exercises no judgment. This runs at machine speed.
%   Phase 3 (Reveal + Record): a opens the commitment by revealing
%     (C, r). The witness publishes Ver(Com, C, E_result) to the
%     ledger. The entry is permanent and publicly auditable.
%
%   Key property: Phase 2 is fully algorithmic. The witness is
%   executing a protocol, not making a judgment call. All judgment
%   was resolved at the governance layer when the benchmark was
%   accepted.

% §5.2 Commitment scheme requirements
%   - Binding: Pr[a opens to C' ≠ C] ≤ negl(λ)
%   - Hiding: no PPT adversary can distinguish Com(C, r) from Com(C', r')
%   - Pedersen commitments satisfy both under the discrete-log assumption.
%   - Discussion of computational vs. information-theoretic binding/hiding
%     trade-offs (Pedersen is computationally binding, perfectly hiding).

% §5.3 Zero-knowledge verification
%   Can the agent prove a claim holds without revealing the claim itself?
%   ZK proof: the agent proves "I satisfy predicate P defined by
%   accepted benchmark B" without revealing which specific claim
%   satisfies it.
%   Formal construction: the agent commits to a claim, the algorithmic
%   witness verifies a predicate on the claim's properties without
%   learning the claim's identity.

% §5.4 Continuous behavioral monitoring
%   Between explicit commitment-verify-reveal cycles, algorithmic
%   witnesses run continuous statistical monitoring:
%   - Is agent a's observed behavior consistent with its last committed
%     claims?
%   - Has the residual between committed baseline and observed behavior
%     drifted beyond the SPRT threshold?
%   This is the "always on" layer that catches behavioral divergence
%   between explicit verification events, operating at machine speed
%   on cross-substrate infrastructure.

% §5.5 Coalition resistance
%   Under substrate exclusivity, a coalition of size k on substrate s
%   cannot influence the verification outcome of an algorithmic witness
%   on substrate s' ≠ s, because the witness executes a fixed protocol
%   without discretion. The coalition's only attack surface is the
%   governance layer (proposing biased benchmarks), which is defended
%   by the fork mechanism of §4.2.

% TODO: formal definitions, protocol specification, security proofs.
